\paragraph{Coordinate expansions and their signed clauses.}
The cardinality test in
Theorem~\ref{thm:damp-fixed-shadow-signed-clutter} parametrizes all ample
realizations at once.  Fixing a contraction and a reduction gives
a complementary constructive description: the remaining choices are
independent labels on connected components.  The following theorem identifies
those labels with the new bits of the signed clauses.  Its conclusion concerns ampleness. In the USO formulation, membership in
$\Damp$ asks whether the constructed support admits any acyclic USO;
forbiddenness additionally asks for independent acyclic USOs on all proper
pc-minors. A labeling alone supplies neither orientation nor those witnesses.
Fixing a restriction and reduction representation map is a more constrained
extension problem than choosing the support labels. The completion theory
in the \href{research/ample_filtrations.pdf}{representation-map writeup}
treats that additional prescribed data. Its fixed-map criteria do not
replace the existential conditions in the present support classification.

\begin{theorem}[Component labels and expansion signs]
\label{thm:damp-component-expansion-signs}
Let $\varnothing\ne C\subseteq A\subseteq Q_F$ be ample, let $e\notin F$,
and let $K_1,\ldots,K_k$ be the connected components of the graph induced
by $A\setminus C$.  For $b=(b_1,\ldots,b_k)\in\{0,1\}^k$, put
\[
 X_b=(C\times\{0,1\})\ \cup\
       \bigcup_{i=1}^k(K_i\times\{b_i\}).
\]
Then $X_b$ is ample, and
\begin{equation}
 \operatorname{Sh}(X_b)=\operatorname{Sh}(A)
       \cup\{T\cup\{e\}:T\in\operatorname{Sh}(C)\}.
 \label{eq:damp-component-expansion-shadow}
\end{equation}
Conversely, these are exactly the ample families on $F\cup\{e\}$ whose
contraction in $e$ is $A$ and whose reduction in $e$ is $C$.
The labels are recovered uniquely.

Write $\mathcal I=\mathcal N(\operatorname{Sh}(C))\cap
\operatorname{Sh}(A)$.  For $T\in\mathcal I$, let $\tau_T\in Q_T$
be the unique trace missing from $C|_T$.  The nonempty fiber
\[
 A(T,\tau_T)=\{x\in A:x|_T=\tau_T\}
\]
lies in a unique component $K_{i(T)}$.  The signed minimal nonfaces of
$X_b$ are exactly
\begin{enumerate}
 \item the signed minimal nonfaces of $A$, on supports not containing $e$;
 \item for each $T\in\mathcal I$, the support $T\cup\{e\}$ with missing
       trace $(\tau_T,1-b_{i(T)})$.
\end{enumerate}
Every component occurs as some $i(T)$.  Thus, with these supports, old
clauses, and traces $\tau_T$ fixed, a choice of new-coordinate bits is
admissible for the shadow in \eqref{eq:damp-component-expansion-shadow}
if and only if its bits agree whenever $i(T)=i(U)$.  There are exactly
$2^k$ such choices.  If $A=C$, the component list and $\mathcal I$ are
empty and the construction is $A\times\{0,1\}$.
\end{theorem}

\begin{proof}
We first prove ampleness directly by comparing shattered and strongly
shattered supports.  We use that complements and face sections of ample
families are ample, and that nonempty ample families are connected
\cite[Theorem~2 and Lemma~1]{Lawrence83}.

Let $Q$ be any cube contained in $A$.  The family $C\cap Q$ is ample
inside $Q$, so its complement $Q\setminus C$, when nonempty, is connected.
Consequently it belongs to one $K_i$.  The entire cube $Q$ therefore
occurs in at least one layer of $X_b$.  A support $T\subseteq F$ is
shattered by $X_b$ exactly when it is shattered by its contraction $A$;
since $A$ is ample, the preceding cube argument shows that it is also
strongly shattered by $X_b$.

A support $T\cup\{e\}$ is strongly shattered by $X_b$ exactly when $T$
is strongly shattered by $C$.  If $T$ is shattered by $C$, both layers
already realize every $T$-trace.  Conversely, suppose $T$ is not shattered
by $C$.  If it is not shattered by $A$, it cannot give a shattered support
of $X_b$.  Otherwise choose a trace $q$ missing from $C|_T$.  The nonempty
face section $A(T,q)$ is connected and avoids $C$, so it lies in one
component $K_i$.  Only layer $b_i$ realizes $q$; hence $T\cup\{e\}$ is
not shattered.  This proves \eqref{eq:damp-component-expansion-shadow}
and equality with the strongly shattered supports, proving ampleness.

For the converse, let $X$ have the stated contraction and reduction.  Every member of $C$ has both lifts, and every member of
$A\setminus C$ has a unique lift.  Two adjacent vertices of
$A\setminus C$ cannot have different lift values: their two lifts would
form an antipodal pair in a square face, with the other two vertices
absent.  Such a face section is not ample.  Thus lift values are constant
on each $K_i$, which gives $X=X_b$ and uniqueness.

The minimal nonfaces of the shadow in
\eqref{eq:damp-component-expansion-shadow} which avoid $e$ are exactly
those of $\operatorname{Sh}(A)$.  A minimal nonface containing $e$ has
the form $T\cup\{e\}$, where $T$ is shattered by $A$, is not shattered
by $C$, and every proper subset of $T$ is shattered by $C$.  These are
precisely the members of $\mathcal I$.  The unique missing trace
$\tau_T$ exists by Theorem~\ref{thm:damp-fixed-shadow-signed-clutter}.
Its fiber in $A$ is a nonempty connected face section outside $C$.
Only its component's layer realizes that trace, giving the stated sign.

To see that every component occurs, take $x\in A\setminus C$.
The signed-clause presentation of $C$ supplies a minimal nonface $T$ of
$\operatorname{Sh}(C)$ with $x|_T=\tau_T$.  All proper subsets of $T$
are shattered by both $C$ and $A$.  If $T$ were not shattered by $A$,
the two projections $C|_T\subseteq A|_T$ would both have
$2^{|T|}-1$ members.  They would be equal, contradicting the presence
of $x|_T$ in $A|_T$.  Hence $T\in\mathcal I$ and $x\in A(T,\tau_T)$.

Finally, bits constant on each component's clause group produce $X_b$
and are admissible by the construction.  Conversely, any admissible
signing with the displayed fixed data defines an ample family $Y$ with
that shadow.  The old clauses imply that its contraction $A'$ is
contained in $A$.  Every point in its reduction avoids each
$(T,\tau_T)$ for $T\in\mathcal I$, because otherwise one of its two
lifts is forbidden.  Every minimal nonface of $\operatorname{Sh}(C)$
outside $\mathcal I$ is a minimal nonface of $\operatorname{Sh}(A)$;
the two projections have the same unique missing trace, by the preceding
cardinality argument.  Thus the old clauses and these additional
conditions imply all clauses defining $C$, so $C'\subseteq C$.  Since
\[
 |Y|=|A'|+|C'|=|A|+|C|,
\]
both inclusions are equalities.  The converse already proved identifies
$Y=X_b$ and forces equality of the bits within each clause group.
\end{proof}

The theorem gives an intrinsic recursion for ample families: choose nested
ample families on fewer coordinates and label their complementary components.
Fixed coordinates may be added separately.  Every active coordinate of a
nonempty ample family has a nonempty reduction, by connectedness,
so the converse supplies exhaustive recovery.  This construction of the
ambient ample class does not identify the forbidden labelings.

\begin{corollary}[Constructing realizations in a fixed coordinate order]
\label{cor:damp-component-reflection-generation}
Fix a simplicial complex $\Sigma\subseteq2^F$ and an ordering
$e_1,\ldots,e_n$ of $F$.  Start with the down-set
\[
 D_\Sigma=\{x\in Q_F:\{f:x_f=1\}\in\Sigma\}.
\]
At stage $j$, form the contraction $A$ and reduction $C$ in
$e_j$ of the current family.  Its unique lifts outside $C$ are all in
layer zero.  Independently choose a label for every component of
$A\setminus C$ and apply the construction of
Theorem~\ref{thm:damp-component-expansion-signs}.  If $C$ is empty,
$e_j$ is constant and the choice is simply to retain or reflect that
constant coordinate.  After these $n$ stages, the possible outputs are
exactly all ample families with shadow $\Sigma$.

For the prescribed coordinate order, every output determines its entire
sequence of intermediate families and component choices uniquely.  In
particular, the graph joining two realizations when one component is
reflected in one coordinate is connected; reflection of a constant
coordinate is included as a move.
\end{corollary}

\begin{proof}
Define $D_e(H)$ by moving every unique lift in coordinate $e$ to layer
zero and retaining every double lift.  The theorem shows that $D_e(H)$
is ample with the same shadow and that its signed-clause presentation
is obtained by setting the $e$-bit to one in every clause containing $e$.
Other bits are unchanged.  For a constant coordinate the same assertion
holds for its singleton clause, and the operation is a coordinate
reflection or the identity.

Consequently the $D_e$ commute.  Applying every one of them makes every
clause the all-one trace, which gives exactly $D_\Sigma$.  This is the
signed-clause form of simultaneous shift normalization.  A family for
which all clause bits in coordinate $e$ are one is downward closed in
$e$: changing that coordinate from one to zero cannot create a forbidden
trace.  Thus all unique lifts at a stage involving an as-yet unprocessed
coordinate lie in layer zero, as asserted.

Reversing the normalization steps, in the prescribed order, reconstructs
any target realization by the component choices in the theorem.  Each
individual component may be reflected separately; all intervening families
are ample with the same shadow.  Conversely, every choice allowed by the
recipe preserves those two properties.  If $H$ is the final output,
the intermediate family after stage $j$ is
\[
 D_{e_{j+1}}\cdots D_{e_n}(H).
\]
Its signed clauses retain exactly the final bits in processed coordinates
and have all-one bits in the remaining coordinates.  These intermediate
families are therefore uniquely determined, as are their recovered
component labels.  Reversing a path to $D_\Sigma$ proves connectivity.
\end{proof}

\begin{proposition}[Alternative repairs in component labelings]
\label{prop:damp-component-label-repair-disjunction}
Let $C\subseteq S\subseteq Q_F$ be nonempty ample families, with $C$
corner-peelable and $S$ not corner-peelable.  Choose distinct points
$a_1,\ldots,a_r$ outside $S$, where $r\ge1$, and put
$R=\{a_1,\ldots,a_r\}$.  Assume:
\begin{enumerate}
 \item $A=S\cup R$ is ample;
 \item each $\{a_i\}$ is a connected component of $A\setminus C$;
 \item every $S\cup\{a_i\}$ is corner-peelable.
\end{enumerate}
In the expansion of Theorem~\ref{thm:damp-component-expansion-signs}, give
all components contained in $S\setminus C$ label zero and give
$\{a_i\}$ label $\varepsilon_i$.  Then
\[
 X_{\varepsilon}\in\Damp
 \quad\Longleftrightarrow\quad
 \varepsilon_i=0\text{ for at least one }i.
\]
If all $\varepsilon_i=1$, the nonpeelable family $S$ is a proper section
of $X_{\varepsilon}$, so this negative expansion is not a forbidden minor.
\end{proposition}

\begin{proof}
Every family $S\cup U$, with $U\subseteq R$, is a component interpolant
between $C$ and $A$.  It is ample: assign its selected components one
label and the others the opposite label in
Theorem~\ref{thm:damp-component-expansion-signs}, and take the appropriate
section.  Thus deleting any available $a_i$ preserves ampleness and is a
corner deletion.  If $U\ne\varnothing$, delete all but one of its members
and finish with the given peeling of $S\cup\{a_i\}$.  In particular,
$A$ and every such $S\cup U$ are corner-peelable.

Suppose some $\varepsilon_i=0$, and let $U$ consist of these tips.
Delete the unique lifts of the tips with label one.  At each step the
remaining family is an ample expansion over a component interpolant
$S\cup U'$ and the same $C$, by the theorem, so each deletion is a
corner deletion.  The remainder is the peripheral expansion with base
$S\cup U$ and site $C$.  Both peel, and
Proposition~\ref{prop:damp-exact-expansion-product-tests} completes its
peeling.  If instead all labels are one, the zero section is exactly
$S$; minor closure of corner peelability proves the negative assertion
and also excludes pc-minimality.
\end{proof}

\begin{remark}[Ample sign choices and peelable sign choices differ]
\label{ex:damp-component-signs-nonlinear-peeling}
The fixed 291-concept cornerless family $S$ has two independent repairing
tips $91$ and $157$, with incident supports $102$ and $390$ in binary
encoding.  Let $C$ be the union of their punctured incident four-cubes.
The fixed replay verifies that $|C|=26$, that $C$ peels, and that
$A=S\cup\{91,157\}$ is ample and peels.  Its complement of $C$ has
components of orders $265,1,1$, the last two being the tips.

There are twelve new-coordinate clauses: ten belong to the first component
and one to each tip.  All eight component labelings give ample expansions
of order $319$.  Precisely six peel.  With the first component labelled
zero, the labelings $(0,1,0)$ and $(0,0,1)$ peel, whereas their
coordinatewise maximum $(0,1,1)$ does not.  Thus the equality constraints
describing admissible expansion signs do not describe the peelable subset.
The negative examples have the proper nonpeelable section $S$, as the
proposition requires; they do not test a criterion restricted to families
whose proper minors all peel.

The independent checker \path{verify_component_expansion_signs.py} also
tests all $2244$ nested nonempty ample pairs through $Q_3$ and all $9656$
new-clause bit assignments.  Exactly $5384$ assignments satisfy the
component equalities and the asserted shadow formula.  It also normalizes
these small families in every coordinate order and replays the inverse
component reflections.  The fixed larger
example is checked by full positive orders and a cornerless proper-section
certificate, without a peeling search for negative answers.  Inputs,
orders, signs and scope are in
\path{component_expansion_signs_verification.json}.  These finite checks
support the examples and the implementation; the universal claims use
the preceding proofs.
\end{remark}
